\chapter[Conclusões]{Conclusões}

O sistema CA é mais flexível e independente de restrições de ordem numérica e geométrica, fato que o coloca em vantagem em relação à solução tradicional da MGDO, na modelagem da propagação de ondas em sistemas físicos. É possível destacar as principais vantagens: 

\begin{itemize}
    \item possibilidade de se trabalhar com células desligadas. A MGDO, bem como outros métodos baseados em FDTD, não permite tal condição.
	\item fácil alteração do comportamento celular, bastando para isso modificar regras ou sinais de controle. A resolução numérica de uma EDO requer condições mais rígidas para convergência e estabilidade, para que se tenha coerência nos resultados obtidos,
	\item apresenta melhor desempenho computacional por estar diretamente envolvido com a discretização do fenômeno emulado. Em um processo envolvendo a resolução de EDO, temos geralmente uma quantidade significativa de cálculos auxiliares, necessários para aproximações numéricas, procedimentos de otimização, dentre outros. 
\end{itemize}

Por outro lado, a formulação de um sistema CA pode exigir um grau mais elevado de criatividade e observação, para que se obtenha boa consistência na emulação do fenômenos físicos explorados. O engenheiro/projetista deve ser capaz de escolher, para cada situação, a estratégia que apresente o melhor equilíbrio entre o grau de realismo requerido e o custo da solução.

Todos os objetivos almejados para este trabalho foram atingidos. Na primeira parte do trabalho, o sistema CA Camphs1D foi desenvolvido, considerando duas células autômatas, com funcionamentos distintos, denominadas GCA1D e HCA1D. Estas células são capazes de operar todos os fenômenos básicos da onda, reproduzindo o comportamento de uma MGDO padrão. Códigos geradores de texto em VHDL, bem como funções decodificadoras de valores binários, foram produzidos para suportar as validações numéricas, que foram realizadas com sucesso.

Com relação às implementações em placa, mais especificamente sobre os sistemas CA SoC FPGA desenvolvidos, este trabalho traz algumas contribuições paralelas. Foi possível verificar a possibilidade de uma aceleração em hardware, realizada pelo PS através do uso de uma interface do tipo \textit{AXI4-Lite}. Também foi possível integrar a este sistema uma interface do tipo \textit{AXI4-Stream} com acesso direto a memória, permitindo realizar um fluxo de dados de entrada, paralelamente à execução dos Co-processadores.

No que se diz respeito somente às interfaces de dados, é possível levantar vantagens e desvantagens relacionadas às suas utilizações. A interface \textit{AXI4-Lite} requer menor quantidade de recursos e, portanto, possibilita a implementação de um CORE CA maior. Por outro lado, esta interface é de baixa velocidade (em relação ao \textit{AXI4-Stream}), possui um menor volume de dados e só permite que o estado final das CA seja externalizado.

A interface \textit{AXI4-Stream}, em contraste com a interface \textit{AXI4-Lite}, oferece uma maior quantidade de dados, com maior velocidade de comunicação. Temos a possibilidade de externalizar sinais (receptores) paralelamente à execução do CORE CA. Por outro lado, a implementação desta interface requer maior quantidade de recursos, e esta restrição acarreta na implementação de um CORE CA reduzido, em relação ao caso anterior.

A estratégia utilizada em todo este trabalho, no que se diz repeito ao desenvolvimento das arquiteturas, contemplou a implementação de um sistema CA completo, em placa. Esta estratégia não se mostrou satisfatória, quando pensada em termos do consumo de recursos implementados. A justificativa, para esta conclusão, está na baixa quantidade de células que podemos implementar no circuito, trabalhando em paralelo - a modelagem de um sistema físico, visando uma aplicação no mundo real, pode necessitar uma quantidade milhares de vezes maior do que as quantidades praticadas neste trabalho.

Ainda que as células autômatas sejam componentes síncronos, os organismos artificiais produzidos podem ser considerados assíncronos, pois os mesmos oferecem todas as condições necessárias para as células operarem em função dos sinais propagados. Deste modo, o sinal de relógio, praticado nas implementações, dita a velocidade de operação do sistema CA, e a latência de suas células autômatas (em ciclos de relógio).

Pensando na transposição deste problema, visando a escalabilidade para aplicações maiores, foi desenvolvida uma solução alternativa para o sistema CA Camphs1D. A denominada estratégia CA fragmentada está apresentada no apêndice I, desenvolvida em nível de software, por representar uma prototipagem inicial. Como parte dos trabalhos futuros, a implementação, em chip, da estratégia CA fragmentada foi desdobrada para um outro trabalho de conclusão de curso (em andamento). Tais implementações devem, obrigatoriamente, ser obtidas em projetos SoC FPGA, devido os benefícios do uso do acesso a memórias do tipo RAM - tal como foi implementado para os organismos artificiais Camphs1D05G1AL e Camphs1D05G1AS.   

%A combinação destes dois barramentos, de dados ao longo de uma simulação, tem potencial de apresentar uma solução mais robusta, principalmente se forem utilizados recursos para a reconfiguração dinâmica do PL. Desta forma, seria possível estruturar um sistema CA que se modifique, em função das suas necessidades, para trabalhar de maneira mais otimizada, somando assim as vantagens dos dois barramentos mencionados. Esta aplicação é parte dos trabalhos futuros, por este motivo não foi explorada neste trabalho, que representa o primeiro nível de desenvolvimento do sistema CA proposto para a modelagem de ondas acústicas em sistemas físicos.